\documentclass{beamer}

\usetheme{Madrid}
\usecolortheme{default}

\usepackage[utf8]{inputenc}
\usepackage[T1]{fontenc}
\usepackage[brazil]{babel}
\usepackage{amsmath, amsfonts, amssymb}

\title{Exemplos de Aplicações de Desigualdades de Concentração}
\subtitle{Large Deviations e Bin Packing}
\author{Com base em Grimmett \& Stirzaker (2020)}
\date{}

\begin{document}

% -----------------------------------------------------
\begin{frame}
  \titlepage
\end{frame}

% -----------------------------------------------------
\begin{frame}{Exemplo (8): Grandes Desvios}
Considere variáveis aleatórias independentes
\[
X_1, X_2,\ldots, X_n \sim \mathrm{Bernoulli}(p).
\]

Defina:
\[
S_n = X_1 + \cdots + X_n, 
\qquad 
Y_n = S_n - np.
\]

Então $(Y_n)$ é um martingal.

\medskip

Pela desigualdade de Hoeffding, para $x>0$:
\[
\mathbb{P}\!\left(|S_n - np| \ge x\sqrt{n}\right)
\le 
2 \exp\!\left( -\frac{x^2}{2\mu^2} \right),
\]
onde 
\[
\mu = \max\{p,\,1-p\}.
\]
\end{frame}

% -----------------------------------------------------
\begin{frame}{Grandes Desvios: Interpretação}
\begin{itemize}
    \item A soma $S_n$ é fortemente concentrada em torno de $np$.
    \item Desvios maiores que ordem $\sqrt{n}$ têm probabilidade exponencialmente pequena.
    \item O resultado é análogo às desigualdades de Bernstein e Chernoff.
\end{itemize}
\end{frame}

% -----------------------------------------------------
\begin{frame}{Exemplo (9): Bin Packing}
Problema clássico: determinar o número mínimo de bins de capacidade 1 necessários para empacotar $n$ objetos.

\medskip

Considere variáveis independentes:
\[
X_1, \ldots, X_n \in [0,1].
\]

Seja:
\[
B_n = \text{número mínimo de bins para empacotar } X_1,\ldots,X_n.
\]

\medskip

\textbf{Fato importante}: existe $\beta > 0$ tal que
\[
\frac{B_n}{n} \longrightarrow \beta
\quad \text{q.c. e em } L^2.
\]
\end{frame}

% -----------------------------------------------------
\begin{frame}{Construção do Martingal}
Para $i \le n$, defina:
\[
Y_i = \mathbb{E}(B_n \mid \mathcal{F}_i),
\qquad
\mathcal{F}_i = \sigma(X_1,\ldots,X_i).
\]

Então:
\[
Y_0 = \mathbb{E}(B_n),
\qquad
Y_n = B_n,
\]
e $(Y_i)$ é um martingal.

\medskip

Além disso, vale:
\[
|Y_i - Y_{i-1}| \le 1.
\]

Logo, Hoeffding se aplica diretamente.
\end{frame}

% -----------------------------------------------------
\begin{frame}{Aplicando Hoeffding ao Bin Packing}
Segue da desigualdade de Hoeffding:
\[
\mathbb{P}\left(|B_n - \mathbb{E}(B_n)| \ge x \right)
\le 2\exp\!\left( -\frac{x^2}{2n} \right),
\qquad x>0.
\]

Em particular, tomando $x = \varepsilon n$:
\[
\mathbb{P}\!\left(|B_n - \mathbb{E}(B_n)| \ge \varepsilon n\right)
\le 
2 \exp\!\left(-\tfrac{1}{2}\varepsilon^2 n\right).
\]

\medskip

Convergência exponencialmente rápida.
\end{frame}

% -----------------------------------------------------
\begin{frame}{Concentração em torno de \texorpdfstring{$\beta n$}{βn}}
Como $B_n/n \to \beta$, temos:
\[
\mathbb{P}\!\left(|B_n - \beta n| \ge \varepsilon n \right)
\le 
2 \exp\!\left(-\tfrac{1}{2}\varepsilon^2 n[1+o(1)]\right).
\]

\bigskip

\textbf{Conclusão:}  
O número ótimo de bins cresce linearmente e apresenta forte concentração em torno de $\beta n$.
\end{frame}

% -----------------------------------------------------
\begin{frame}{Conclusões}
\begin{itemize}
    \item Hoeffding é uma ferramenta essencial para controle de flutuações.
    \item No exemplo de Bernoulli, controla grandes desvios da soma.
    \item No bin packing, mostra que o número de bins raramente foge muito do valor médio.
    \item Estratégia comum: construir martingais com incrementos limitados.
\end{itemize}
\end{frame}

\end{document}
\documentclass{beamer}

\usetheme{Madrid}
\usecolortheme{default}

\usepackage[utf8]{inputenc}
\usepackage[brazil]{babel}
\usepackage{amsmath, amsfonts, amssymb}

\title{Exemplos de Aplicações de Desigualdades de Concentração}
\subtitle{Large Deviations e Bin Packing}
\author{Com base em Grimmett \& Stirzaker (2020)}
\date{}

\begin{document}

\frame{\titlepage}

%--------------------------------------------------
\begin{frame}{Exemplo (8): Grandes Desvios}
Considere variáveis aleatórias independentes 
\[
X_1, X_2, \ldots, X_n \sim \text{Bernoulli}(p).
\]

Seja
\[
S_n = X_1 + \cdots + X_n, 
\qquad
Y_n = S_n - np.
\]

\medskip

Então $(Y_n)$ é um martingal, e pela desigualdade de Hoeffding vale, para $x>0$:
\[
\mathbb{P}\left(|S_n - np| \geq x\sqrt{n}\right)
\le 
2 \exp\left( -\frac{1}{2}\frac{x^2}{\mu^2} \right),
\]
onde 
\[
\mu = \max\{p,\;1-p\}.
\]

\medskip

\textbf{Interpretação:} As flutuações de $S_n$ em torno de sua média $np$ têm probabilidade exponencialmente pequena quando ultrapassam uma escala da ordem de $\sqrt{n}$.
\end{frame}

%--------------------------------------------------
\begin{frame}{Grandes Desvios: Comentários}
\begin{itemize}
    \item A desigualdade é análoga às desigualdades de Bernstein e Chernoff.
    \item Captura o fenômeno típico de \textbf{grandes desvios}: a soma de variáveis Bernoulli raramente se desvia muito de sua média.
    \item Resultado fundamental em estatística, aprendizado de máquina e teoria da informação.
\end{itemize}
\end{frame}

%--------------------------------------------------
\begin{frame}{Exemplo (9): Bin Packing}
Problema clássico de pesquisa operacional:  
\textbf{“Quantos bins (caixas) de capacidade 1 são necessários para acomodar $n$ objetos com tamanhos aleatórios?”}

\medskip

Considere:
\[
X_1, X_2, \ldots, X_n \text{ independentes com distribuição em } [0,1].
\]

Defina $B_n =$ número mínimo de bins necessários para empacotar todos os objetos eficientemente.

\medskip

\textbf{Fato importante:} Existe $\beta>0$ tal que
\[
\frac{1}{n} B_n \longrightarrow \beta
\quad \text{quase certamente e em média quadrática}.
\]
\end{frame}

%--------------------------------------------------
\begin{frame}{Bin Packing: A Construção do Martingal}
Para $i\le n$, defina
\[
Y_i = \mathbb{E}(B_n \mid \mathcal{F}_i),
\]
onde $\mathcal{F}_i = \sigma(X_1,\ldots,X_i)$.

Então $(Y_i)$ é um martingal, com:
\[
Y_0 = \mathbb{E}(B_n),
\qquad
Y_n = B_n.
\]

\medskip

Um argumento sobre adicionar ou remover um item mostra que:
\[
|Y_i - Y_{i-1}| \le 1.
\]

Assim, Hoeffding se aplica.
\end{frame}

%--------------------------------------------------
\begin{frame}{Bin Packing: Aplicando Hoeffding}
A desigualdade de Hoeffding fornece:
\[
\mathbb{P}\left( |B_n - \mathbb{E}(B_n)| \ge x \right)
\le 
2 \exp\left( -\frac{x^2}{2n} \right), 
\qquad x>0.
\]

Em particular, para $x = \varepsilon n$:
\[
\mathbb{P}\left( |B_n - \mathbb{E}(B_n)| \ge \varepsilon n \right)
\le 
2\exp\left( -\frac{1}{2}\varepsilon^2 n \right),
\]
decadência exponencial.
\end{frame}

%--------------------------------------------------
\begin{frame}{Bin Packing: Convergência em Probabilidade}
Como $B_n/n \to \beta$, obtemos:
\[
\mathbb{P}\left(|B_n - \beta n| \ge \varepsilon n\right)
\le
2\exp\left(-\tfrac{1}{2}\varepsilon^2 n[1+o(1)]\right).
\]

\bigskip

\textbf{Conclusão:}  
$B_n$ se concentra fortemente em torno de $\beta n$.
\end{frame}

%--------------------------------------------------
\begin{frame}{Conclusões da Aula}
\begin{itemize}
    \item A desigualdade de Hoeffding é uma ferramenta fundamental para estudar concentração.
    \item Em (8), controla grandes desvios da soma de Bernoullis.
    \item Em (9), mostra que o número ótimo de bins oscila pouco em torno de seu valor médio.
    \item Ideia central: \textbf{construir martingais} cujo incremento é uniformemente limitado.
\end{itemize}
\end{frame}

\end{document}
